\begin{thebibliography}{99}
\footnotesize\itemsep=0pt

\bibitem{And67}
Andrews G.E., On {S}chur's second partition theorem,
 \href{https://doi.org/10.1017/S0017089500000197}{\textit{Glasgow Math.~J.}}
 \textbf{8} (1967), 127--132.

\bibitem{Car10}
Carter R.W., Lie algebras of finite and affine type, \textit{Cambridge Stud.
 Adv. Math.}, Vol.~96,
 \href{https://doi.org/10.1017/CBO9780511614910}{Cambridge University Press},
 Cambridge, 2005.

\bibitem{D17}
Dousse J., The method of weighted words revisited, \textit{S\'em. Lothar.
 Combin.} \textbf{78B} (2017), 66, 12~pages,
 \href{http://arxiv.org/abs/1702.07177}{arXiv:1702.07177}.

\bibitem{DHK25}
Dousse J., Hardiman L., Konan I., Partition identities from higher level
 crystals of~{$\smash{A_1^{(1)}}$},
 \href{https://doi.org/10.1090/proc/16417}{\textit{Proc. Amer. Math. Soc.}}
 \textbf{153} (2025), 1363--1382,
 \href{http://arxiv.org/abs/2111.10279}{arXiv:2111.10279}.


\bibitem{DK22a}
Dousse J., Konan I., Generalisations of {C}apparelli's and {P}rimc's
 identities,~{I}: {C}oloured {F}robenius partitions and combinatorial proofs,
 \href{https://doi.org/10.1016/j.aim.2022.108571}{\textit{Adv. Math.}}
 \textbf{408} (2022), 108571, 70~pages,
 \href{http://arxiv.org/abs/1911.13191}{arXiv:1911.13191}.

\bibitem{DK19b}
Dousse J., Konan I., Generalisations of Capparelli's and Primc's identities,
 {II}: {P}erfect $\smash{A_{n-1}^{(1)}}$ crystals and explicit character
 formulas, \href{http://arxiv.org/abs/1911.13189}{arXiv:1911.13189}.

\bibitem{DK22}
Dousse J., Konan I., Multi-grounded partitions and character formulas,
 \href{https://doi.org/10.1016/j.aim.2022.108275}{\textit{Adv. Math.}}
 \textbf{400} (2022), 108275, 41~pages,
 \href{http://arxiv.org/abs/2103.04983}{arXiv:2103.04983}.

\bibitem{DK22c}
Dousse J., Konan I., Characters of level $1$ standard modules of
 $\smash{C_n^{(1)}}$ as~generating functions for generalised partitions,
 \href{http://arxiv.org/abs/2212.12728}{arXiv:2212.12728}.

\bibitem{DL18}
Dousse J., Lovejoy J., On a {R}ogers--{R}amanujan type identity from crystal
 base theory, \href{https://doi.org/10.1090/proc/13694}{\textit{Proc. Amer.
 Math. Soc.}} \textbf{146} (2018), 55--67,
 \href{http://arxiv.org/abs/1612.05423}{arXiv:1612.05423}.

\bibitem{HK02}
Hong J., Kang S.-J., Introduction to quantum groups and crystal bases,
 \textit{Grad. Stud. Math.}, Vol.~42,
 \href{https://doi.org/10.1090/gsm/042}{American Mathematical Society},
 Providence, RI, 2002.

\bibitem{Kac90}
Kac V.G., Infinite-dimensional {L}ie algebras, 3rd ed.,
 \href{https://doi.org/10.1017/CBO9780511626234}{Cambridge University Press},
 Cambridge, 1990.

\bibitem{KMN92}
Kang S.-J., Kashiwara M., Misra K.C., Miwa T., Nakashima T., Nakayashiki A.,
 Perfect crystals of quantum affine {L}ie algebras,
 \href{https://doi.org/10.1215/S0012-7094-92-06821-9}{\textit{Duke Math.~J.}}
 \textbf{68} (1992), 499--607.

\bibitem{KMOY07}
Kashiwara M., Misra K.C., Okado M., Yamada D., Perfect crystals
 for~{$\smash{U_q\bigl(D^{(3)}_4\bigr)}$},
 \href{https://doi.org/10.1016/j.jalgebra.2007.02.021}{\textit{J.~Algebra}}
 \textbf{317} (2007), 392--423,
 \href{http://arxiv.org/abs/math.QA/0610873}{arXiv:math.QA/0610873}.

\bibitem{Lep78}
Lepowsky J., Lectures on Kac--Moody Lie algebras, Universit\'e Paris~VI, Paris,
 1978.

\bibitem{LM78}
Lepowsky J., Milne S., Lie algebraic approaches to classical partition
 identities, \href{https://doi.org/10.1016/0001-8708(78)90004-X}{\textit{Adv.
 Math.}} \textbf{29} (1978), 15--59.

\bibitem{LW84}
Lepowsky J., Wilson R.L., The structure of standard modules.~{I}. {U}niversal
 algebras and the {R}ogers--{R}amanujan identities,
 \href{https://doi.org/10.1007/BF01388447}{\textit{Invent. Math.}} \textbf{77}
 (1984), 199--290.

\bibitem{LW85}
Lepowsky J., Wilson R.L., The structure of standard modules.~{II}. {T}he
 case~{$\smash{A^{(1)}_1}$}, principal gradation,
 \href{https://doi.org/10.1007/BF01388515}{\textit{Invent. Math.}} \textbf{79}
 (1985), 417--442.

\bibitem{Mac16}
MacMahon P.A., Combinatory analysis. {V}ol.~{II}, Cambridge University Press,
 Cambridge, 1916.

\bibitem{RR19}
Rogers L.J., Ramanujan S., Proof of certain identities in combinatory analysis,
 \textit{Proc. Cambridge Philos. Soc.} \textbf{19} (1919), 211--216.

\bibitem{Sch17}
Schur I., Ein {B}eitrag zur additiven {Z}ahlentheorie und zur {T}heorie der
 {K}ettenbr\"uche, \textit{S.-B.~Preuss. Akad. Wiss. Phys. Math. Klasse}
 (1917), 302--321.

\end{thebibliography}
